\section{模型假设}

为在题面数据范围内建立可闭合模型，将药材视为轴对称均匀圆柱，半径为
$R_0=0.02\,\mathrm m$、长度为 $L=0.25\,\mathrm m$，烘干主体传递沿径向发生。药材内部温度
$T(r,t)$ 和干基含水率 $C(r,t)$ 视为连续场，中心满足零梯度，真实表面满足有效换热和传质
Robin 条件。该一维径向近似适用于研究主体为中截面的情形；端面附近或形状明显偏离圆柱时，
需要用二维轴对称或三维模型复核。

四个子问题均从题设给出的均匀初态开始：
\[
T(r,0)=301.15\,\mathrm K,\qquad C(r,0)=2.55\,\mathrm{kg/kg}.
\]
由于初始含水率与表面平衡条件可能不相容，计算在 $t=0$ 保留题设初值，从 $t>0$ 开始施加表面边界，
避免人为构造没有数据依据的初始边界层。

附件 1 的烘房温度和环境浓度在相邻观测节点之间采用分段线性插值。问题三和问题四超过
$4\,\mathrm h$ 后采用末小时均值延拓；这代表明确的边界情景，而不是对未观测时段的实测恢复。
本文把环境浓度解释为有效表面平衡浓度 $C_e(t)$，并与药材干基含水率严格区分。题面没有给出
水活度或吸附等温线，因此不额外引入无法由现有数据确定的换算参数。

问题一使用附录 2，问题二和问题三使用附录 3，问题四使用附录 4 的整组物性；未给出替代值时
沿用题面换热系数 $h$ 和传质系数 $k_m$。温度方程采用无显式蒸发潜热、无内部热源的有效显热闭合，
水分扩散率的非线性保留在界面通量中。问题四采用附件给定的半径轨迹和均匀径向收缩，假设干骨架无损失，
并在材料坐标 $\xi=r/R(t)$ 中计算后映射回实际半径；守恒核算使用干骨架密度 $\rho_d$，不与热物性密度
$\rho(C)$ 混用。

问题三在固定半径下延续问题二的完整耦合解，问题四在收缩材料坐标下重新求解。两问的达标条件均作用于
连续重构场的全域最大含水率，事件时刻由连续求根确定，而不是由五个展示位置或输出步长直接指定。
内部统一使用 SI 单位，表格和图件仅在输出时换算为 $\mathrm{^\circ C}$、cm 和 h；显示位数不作为真实误差界。

上述假设构成题面信息不足条件下的闭合模型。网格加密、时间收紧、干基水量收支和 Robin 残差用于检验
离散实现的一致性，但不能替代对有效表面平衡关系、经验物性、无显式潜热、端部忽略和均匀收缩等假设的
实验识别。长期结果还需明确 $4\,\mathrm h$ 后环境延拓口径。
